\section{Retour, mémoire et universalité de la cascade dodécaphasique}
\label{sec:hmt-feig-interpretacion}

\subsection{La structure génératrice et la lecture critique}

La construction de Feigenbaum en Holographie Modulaire Triadique relie une
organisation arithmétique orientée à une loi asymptotique de changement d’échelle.
Son point de départ est la composition effective d’APP, TRIT et TPK. APP évalue
la somme et le produit sur les deux feuilles du réseau de chiffres et conserve les
quotients et les restes de ces opérations. TRIT détermine le régime local et
l’orientation de sa lecture. TPK sélectionne et transporte les états,
actualise leurs registres et prolonge leurs routes avec retenue, frontière et mémoire.
La structure discrète conjointe du continu fournit la résidence commune
de ces états et des lecteurs qui agissent sur eux.

Le prolongement
\[
w_6\longrightarrow w_{12}\longrightarrow w_{18}
\longrightarrow w_{24}\longrightarrow w_{30}
\longrightarrow R_{36}\longrightarrow G_9
\]
organise cette persistance. La graine initiale est relevée en conservant ses
préfixes ; les relèvements successifs changent le régime sans perdre l’orientation
ni la mémoire ; la frontière coinductive permet de poursuivre la construction, et les
neuf cartes réunissent ses lectures locales. L’holonomie nonadique exprime le
retour de phase accompagné d’un accroissement de mémoire. Ainsi, un tour
visible conserve sa position dans l’histoire complète. L’opération de retour
acquiert dès l’origine deux composantes inséparables :
réitération d’une règle et conservation de sa provenance.

Le lecteur critique fixé sur la roue dodécaphasique orientée convertit cette
structure en douze phases avec une trace uniforme. Leurs représentants sont
\[
\phi_m=\frac{(3m-1)\pi_{\mathrm{HMT}}}{18},
\qquad w_m=\frac1{12},\qquad 1\le m\le12.
\]
Le choix du représentant orienté importe parce que le deuxième moment
utilise sa grandeur. Le quotient angulaire enregistre la phase, tandis que le
représentant conserve l’information nécessaire à sa normalisation. La
coordonnée \(\pi_{\mathrm{HMT}}\) provient des lecteurs coinductifs du même
noyau APP--TRIT--TPK. Dans son image corrélée, les coordonnées
\(\pi,\varphi,e\) et la clôture dodécaphasique de \(\alpha_{\mathrm{HMT}}\)
conservent leur ordre constructif et leurs domaines propres. Dans le lecteur critique
considéré ici, on utilise la coordonnée angulaire déjà générée.

La normalisation est déterminée par une somme entière :
\[
\sum_{m=1}^{12}(3m-1)^2=5394=6\cdot899,
\qquad k_m=\frac{2(3m-1)}{\sqrt{899}}.
\]
Le facteur commun \(\pi_{\mathrm{HMT}}\) se simplifie lors de la fixation du deuxième moment.
Cette simplification conserve la provenance des phases et produit une
représentation particulièrement précise de leur lecture normalisée. On obtient
\[
u_0(x)=\frac1{12}\sum_{m=1}^{12}\bigl(1-\cos(k_mx)\bigr),
\qquad f_\lambda(x)=1-\lambda u_0(x),
\qquad \sum_mw_mk_m^2=2.
\]
Le profil est pair et entier ; il satisfait \(u_0(0)=u_0'(0)=0\) et
\(u_0''(0)=2\). Pour \(0<\lambda\le2/u_0(1)\), la famille envoie
\([-1,1]\) dans lui-même et possède un maximum critique quadratique avec
\(f_\lambda''(0)=-2\lambda\). Le paramètre \(\lambda\) parcourt cette famille
après fixation du lecteur, de l’orientation, des poids et de l’échelle. Le nombre
universel qui sera reconnu à la fin est ainsi séparé de la sélection
des coefficients.

\subsection{Douze phases, moments et représentation spectrale}

Le profil admet une deuxième expression qui réunit ses phases dans un opérateur
fini. Les nœuds
\[
s_m=k_m^2=\frac{4(3m-1)^2}{899}
\]
déterminent la mesure positive \(\mu_0=12^{-1}\sum_m\delta_{s_m}\). Ses
moments \(M_r=\int s^r\,d\mu_0(s)\) définissent des formes de Hankel strictement
positives jusqu’au rang douze. Cette positivité provient de l’évaluation
d’un polynôme sur douze nœuds distincts : un polynôme de degré inférieur à
douze conserve au moins une évaluation non nulle. Au degré douze, le produit
des douze facteurs nodaux annule la norme et détermine exactement le rang
du lecteur.

L’orthogonalisation de ces moments produit une matrice de Jacobi
\(J_{12}\), réelle symétrique, à coefficients sous-diagonaux positifs. Son calcul fonctionnel
restitue le profil complet :
\[
u_0(x)=e_0^{\mathsf T}
\bigl[I-\cos(x\sqrt{J_{12}})\bigr]e_0.
\]
La somme de cosinus et l’opérateur de Jacobi sont deux réalisations du même
lecteur. Les moments maintiennent toutes les fréquences et leurs poids ; les
marques d’orientation, de route et de mémoire restent dans l’état qui les
engendre. Le rang douze identifie la dimension spectrale de cette lecture
particulière. La profondeur de ses prolongements relève d’une autre dimension
de la construction : la longueur des histoires compatibles et la résolution
avec laquelle elles sont évaluées.

La géométrie analytique est elle aussi explicitée. Pour la déformation
pondérée \(u_\eta\), dans la bande \(|\eta|\le1/40\), les poids restent
positifs et le deuxième moment demeure normalisé. Dans
\(|\operatorname{Re}z|<\pi/3\), il existe une fonction holomorphe impaire et univalente
\(b_\eta\), normalisée par \(b_\eta'(0)=1\), telle que
\(u_\eta=b_\eta^2\). La roue produit donc une déformation conforme
du pli quadratique. Sa conjugaison dynamique conserve cette déformation
à travers \(b_\eta(1-\lambda w^2)\). Le théorème de changement d’échelle paramétrique
développé ici concerne le membre uniforme \(\eta=0\) ; les
propriétés structurelles de la bande pondérée et la persistance locale de
ses racines conservent leurs quantificateurs spécifiques.

\subsection{Profondeur de lecture et restitution des histoires}

La profondeur nonadique affine la lecture d’un paramètre, tandis que le
doublement dynamique change l’horizon de retour. Le raffinement avec
retenue exprime la première opération par
\[
L_c(x,y)=(Bx-c,By+c),\qquad
C_m=\sum_{j=1}^m c_jB^{m-j},
\]
\[
(x_m,y_m)=(B^mx_0-C_m,B^my_0+C_m).
\]
Le support extérieur croît comme \(B^m\), et le cylindre normalisé du registre
a pour diamètre \(B^{-m}\). La paire entière, la profondeur et la marque de
frontière conservent la reconstruction des quotients, des restes et des préfixes.
Croissance extérieure et résolution intérieure décrivent ainsi des opérations
compatibles sur une même histoire.

La composition exacte
\(L_D^{[b]}L_C^{[a]}=L_{B^bC+D}^{[a+b]}\) permet de regrouper une histoire en
blocs sans altérer ses données. Un retour de doublement réunit deux
transitions ; le calendrier nonadique en réunit neuf. Les deux regroupements sont
comparés sur un horizon commun de \(9\,2^N\) événements, en conservant
l’ordre des sous-blocs et leurs retenues.

La contribution de la mémoire peut s’observer directement dans les branches
inverses du profil. Si \(v=(u_0|_{[0,1]})^{-1}\), les deux branches sont
\[
\beta_\sigma(y)=\sigma v\!\left(\frac{1-y}{\lambda}\right),
\qquad \sigma\in\{-1,+1\}.
\]
La valeur finale, jointe à la dernière feuille, reconstruit la valeur précédente. Une
suite de feuilles permet de restituer toutes les étapes intermédiaires. Le mot
des branches accompagne les registres APP--TRIT--TPK de route, de frontière, de retenue et de
mémoire ; chaque registre conserve sa fonction. Cette restitution concrète
explique quelle information exige l’inversion de la lecture critique.

La résolution de chaque retour dispose également d’une borne explicite. Avec
\(t=\lambda u_\eta(1)/2\), \(z=(x+1)/2\) et
\[
F_{t,\eta}(z)=1-t\frac{u_\eta(2z-1)}{u_\eta(1)},
\]
on obtient
\[
\left|F_{t,\eta}^{q}(1/2)-F_{s,\eta}^{q}(1/2)\right|
\le S_q|t-s|,\qquad S_q=\sum_{j=0}^{q-1}(6\sqrt2)^j<9^q.
\]
Un cylindre de base neuf et de profondeur \(m=2^N+r\) permet donc de lire
le retour d’horizon \(2^N\) avec un diamètre inférieur à \(9^{-r}\). Cette
estimation donne un contenu quantitatif à la profondeur arbitraire : pour chaque
horizon et chaque précision demandés, elle détermine une profondeur suffisante.
L’évaluation des fonctions analytiques et l’arrondi possèdent leurs propres
restes, qui se propagent avec l’incertitude du paramètre.

Les vacances interviennent dans le bilan de capacité du registre. La
comparaison entre blocs de 729 et 1000 possibilités distingue la capacité
accumulée de la longueur chronologique, en conservant les étapes où la première
reste constante. Cette comptabilité détermine l’espace requis pour
conserver une lecture et sa profondeur. L’admissibilité des histoires
découle, quant à elle, des règles de survie. Les deux opérations
collaborent à la construction et maintiennent des objets mathématiques distincts.

\subsection{Chronologie nonadique et combinatoire de doublement}

Le mot critique de doublement possède une expression à toute profondeur :
\[
\tau_j=(-1)^{v_2(j)},\qquad
\tau_{2j}=-\tau_j,\qquad \tau_{2j+1}=+1.
\]
Ses préfixes satisfont \(W_{n+1}=W_n(-1)^nW_n\). La règle détermine comment
les deux copies du retour sont orientées et où apparaît leur séparation. La
lecture nonadique de la chronologie conserve l’entier complet
\(K=\rho_9(K)+9q_9(K)\). Lorsqu’on le réduit modulo \(2^n\), l’holonomie qui
avance de neuf événements induit la translation de neuf sur ce cycle. Puisque
neuf est impair, cette translation visite toutes ses positions. La
permutation \(j\mapsto9j\pmod{2^n}\), qui la conjugue à l’avancée unitaire,
conserve la valuation dyadique de chaque position non nulle. Le reste et le quotient
chronologiques soutiennent conjointement cette compatibilité.

La mémoire bilatérale possède également une réalisation opératorielle compatible
avec le raffinement. Sur les cycles de longueur \(2^n\), l’opérateur
d’avancée \(C_n\) est incorporé dans
\[
\mathbb U_n=P_0\otimes I+P_1\otimes C_n,
\]
avec les projecteurs complémentaires du lecteur nonadique. À la limite, ses
composantes visible et complémentaire satisfont
\[
\|T_{\infty,a}v\|^2+\|D_{\infty,a}v\|^2=\|v\|^2,
\qquad
v=T_{\infty,a}^*T_{\infty,a}v+D_{\infty,a}^*D_{\infty,a}v.
\]
Ces identités décrivent la conservation et la restitution dans l’espace
et la représentation déclarés. La contraction des applications renormalisées,
qui apparaîtra ensuite, mesure une autre opération : le rapprochement de leurs profils
vers une forme stable. L’histoire conservée par le relèvement et la
convergence d’une lecture fonctionnelle s’articulent à des niveaux différents.

\subsection{La carte ordonnée et la sélection globale des racines}

La lecture archimédienne du continu HMT transporte les opérations
par un isomorphisme de corps ordonnés complets, dans l’univers
déclaré \(\mathfrak G\) :
\[
\iota:\mathbb R_{\mathrm{HMT}}^{\mathfrak G}
\longrightarrow\mathbb R^{\mathfrak G}.
\]
La construction interne du cosinus par sa série, de la racine positive de 899,
du profil et de ses itérées précède la sélection des zéros. La conservation
de la somme, du produit, de l’ordre et des limites convergentes donne
\[
\iota\bigl(S_n^{\mathrm H}(a)\bigr)=S_n\bigl(\iota(a)\bigr),
\qquad S_n(\lambda)=f_\lambda^{2^n}(0).
\]
La surjectivité couvre toutes les racines du domaine de cette carte ; l’ordre
conserve les retours antérieurs, la période exacte et la condition de
se trouver à droite de la racine précédente. On transporte ainsi
également le minimum de chaque ensemble de candidats. La généalogie atteint
la sélection effective du paramètre, la deuxième projection et ses
marques étant conservées sur le même antécédent.

Le sélecteur est fixé par \(\lambda_1=1/u_0(1)\) et par la première racine
nouvelle de période critique exacte \(2^n\) située à droite de
\(\lambda_{n-1}\). La classification des itinéraires démontre que chaque racine
sélectionnée a pour mot \(W_n0\). Les retours restrictifs réduisent sa
période de moitié en conservant l’ordre unimodal, et l’induction restitue
le mot complet. L’analyticité des retours non identiquement nuls
isole leurs zéros dans chaque compact ; la
comparaison des signes avant le premier paramètre d’itinéraire infini
garantit l’existence d’une racine nouvelle à tous les niveaux.

Soit \(\Lambda_\tau\) l’ensemble compact non vide des paramètres qui réalisent
le mot infini, et soit \(a_*\) son minimum. On démontre
\[
\lambda_n\nearrow a_*.
\]
Le passage à la limite conserve chaque signe au moyen d’une borne uniforme à horizon
fixe. Si \(c_p=f_\lambda^p(0)\), les signes opposés de \(c_p\) et
\(c_{2p}\), avec \((f_\lambda^p)'(0)=0\), impliquent
\[
|c_p|>\frac2{B_p},\qquad
B_p=16^p\frac{16^p-1}{15}.
\]
Ainsi, les préfixes de périodes croissantes préservent le mot infini en
s’accumulant. La preuve utilise la suite initiale des premières racines dès
sa définition globale.

\subsection{La feuille stable et l’identification de la limite sélectionnée}

Le retour normalisé
\[
(Rf)(x)=-\frac{f(f(-ax))}{a},\qquad a=-f(1)>0,
\]
réunit deux transitions, change l’échelle et conserve l’orientation nécessaire
pour restituer un maximum de valeur un. Les domaines restrictifs établissent
son interprétation comme retour réel. La dérivée de cet opérateur inclut
la variation de l’échelle \(a\) ; cette contribution détermine la sensibilité
paramétrique de la transformation complète.

Dans la carte \(f(x)=1-x^2V((x^2-1)/(5/2))\), de coordonnées \(u=10v_0\)
et \(\nu=(v_1,v_2,\ldots)\), le quatrième retour de la famille entre dans un
voisinage quantifié du point fixe \(g\). La courbe traverse transversalement
sa feuille stable en un unique paramètre local
\[
1.874038<\lambda_*<1.874039,
\]
et satisfait
\[
\|R^{4+n}f_{\lambda_*}-g\|_{\mathcal A}
<0.001666(0.83996)^n.
\]
La feuille stable est un graphe de pente au plus \(0.16\). Les
sécantes paramétriques appartiennent à un cône transversal expansif. Les deux
propriétés distinguent la diminution des déformations stables de la
croissance de la séparation entre paramètres.

L’identification de \(a_*\) et de \(\lambda_*\) réunit cette géométrie locale
et le minimum global. On exclut d’abord l’intervalle entre la huitième racine
et l’extrémité inférieure de la boîte locale. Ensuite, l’ordre
\(a_*\le\lambda_*\) permet d’étudier une séparation hypothétique au moyen
du cône négatif. Tant que sa coordonnée transversale a une valeur absolue inférieure
à \(0.006\), la sécante reste dans le domaine analytique et croît d’un
facteur compris entre \(4.4034\) et \(8.0887\). Une séparation persistante
produirait donc une première sortie d’amplitude comprise entre
\(0.006\) et \(0.0485322\).

Le résidu de l’orbite stable conserve une information plus précise que sa
norme : ses composantes satisfont \(|e_u|\le0.16\|e_\nu\|_1\). Cette
relation permet de borner conjointement leurs effets sur l’orbite critique.
Le recouvrement de toute la région de première sortie montre dans chaque cas un
signe incompatible avec \(\tau\), ou un retour contractant qui impose à deux
valeurs critiques dyadiques successives d’avoir le même signe. Le mot
infini exige des signes opposés. La contradiction identifie les paramètres :
\[
\lambda_n\nearrow a_*=\lambda_*.
\]
La démonstration combine une induction valable à toute profondeur avec une
exclusion finie uniforme de la région qu’atteindrait toute
séparation. La portée infinie repose sur cette composition.

\subsection{Changement d’échelle universel et signification de la précision}

La réalisation holomorphe de degré deux des retours et la transversalité
établie permettent de comparer leurs classes hybrides. Pour chaque période
suffisamment élevée, la classe canonique possède une unique intersection locale
avec la famille. Les premières racines globales convergent déjà vers le croisement et
conservent la combinatoire de doublement ; elles coïncident donc à partir d’un certain rang
avec les centres locaux lorsque leurs périodes physiques sont égalées. La preuve
transporte la sélection initiale jusqu’à la loi métrique universelle.

Il existe \(C>0\), \(M<\infty\) et \(0<\theta<1\) tels que
\[
\lambda_*-\lambda_n=C\delta_{\mathrm F}^{-n}(1+e_n),
\qquad |e_n|\le M\theta^n,
\]
et, par conséquent,
\[
\lim_{n\to\infty}
\frac{\lambda_{n-1}-\lambda_{n-2}}
{\lambda_n-\lambda_{n-1}}=\delta_{\mathrm F}.
\]
Le facteur \(\delta_{\mathrm F}\) est la valeur propre expansive de l’opérateur de doublement
reconnu à la fin de la construction. La famille particulière conserve ses
douze phases, ses moments et son histoire de provenance ; la limite de ses
rapports paramétriques appartient à la même classe universelle de criticité
quadratique. Cette relation explique conjointement la spécificité du
générateur et l’universalité de la sortie.

La précision possède ici trois résidences. Les intervalles des racines
contrôlent leur évaluation finie. L’estimation
\(0.001666(0.83996)^n\) contrôle la convergence opératorielle sur la feuille
stable. Le terme \(M\theta^n\) contrôle le transfert de cette géométrie à
l’échelle paramétrique au moyen de l’holonomie transversale. Chaque estimation
correspond à sa variable et à son opération. Le huitième quotient est encadré
autour de \(4.669212189150204154556096215889663928\), avec une largeur inférieure
à \(5.808\cdot10^{-37}\) ; cette largeur mesure l’évaluation de \(\delta_{\mathrm F,8}\).
La distance de \(\delta_{\mathrm F,8}\) à la limite exige de quantifier les constantes
du reste paramétrique, dont l’existence est déjà démontrée.

La notation maintient séparés \(\delta_{\mathrm F}\), facteur de changement d’échelle paramétrique
de Feigenbaum ; \(\alpha_{\mathrm F}\), constante spatiale associée à son
problème de renormalisation ; et \(\alpha_{\mathrm{HMT}}\), constante de
structure fine du corpus. Le présent résultat identifie \(\delta_{\mathrm F}\)
au moyen des premières racines du secteur uniforme \(\eta=0\). La bande
pondérée conserve ses théorèmes structurels et d’existence, et son extension
métrique quantifiée constitue un énoncé spécifique distinct.

L’apport de l’enchaînement réside dans la continuité entre génération,
lecture, mémoire, sélection et limite. L’arithmétique orientée détermine le
profil ; la représentation spectrale conserve ses douze composantes ; les
cylindres permettent de le résoudre avec une précision arbitraire à chaque horizon ;
le transport des histoires restitue ses parcours ; la carte ordonnée
conserve toutes les racines et leurs minima ; la classification prolonge le
sélecteur ; et l’expansion transversale, jointe à la contraction stable,
détermine sa loi asymptotique. L’universalité apparaît comme une propriété
démontrée de cette construction complète et de sa réalisation analytique.
